\section{Setup and Preliminaries}
\label{sec:setup}

We analyze a generic two-module compositional model. While our motivating
application is infrared spectroscopic imaging for tissue classification,
the framework applies to any pipeline that first reduces a per-element
spectral signal and then applies a high-capacity spatial model. We
deliberately keep the notation domain-agnostic.

\subsection{The compositional model}

Let $X \in \R^{H \times W \times S}$ denote a hyperspectral input image,
where $H$ and $W$ are spatial dimensions and $S$ is the spectral
dimension.\footnote{Multi-channel inputs (e.g.\ raw absorbance together
with first and second derivatives) are handled by concatenating channels
along the spectral dimension; this changes only the input size and
leaves the analysis unaffected.} We write the per-pixel spectrum at
location $(h,w)$ as $x_{h,w} \in \R^S$.

The model consists of two modules applied sequentially:
\begin{enumerate}
    \item A \emph{spectral reduction} $f_\thetap: \R^S \to \R^K$
          applied independently to every pixel, producing a feature map
          $Z \in \R^{K \times H \times W}$ with
          $Z_{:,h,w} = f_\thetap(x_{h,w})$.
    \item A \emph{spatial model} $g_\phip: \R^{K \times H \times W} \to
          \R^{N_{\mathrm{cls}} \times H \times W}$ producing per-pixel
          class logits $\hat{y} = g_\phip(Z)$.
\end{enumerate}
We denote the composed model
$F_{\thetap,\phip}(X) = g_\phip(f_\thetap(X))$ and the parameter counts
$\Cf := \dim(\thetap)$ and $\Cg := \dim(\phip)$.

\begin{assumption}[Linear spectral reduction]
\label{ass:linear}
$f_\thetap$ is linear in $\thetap$: there exists $W \in \R^{S \times K}$
with $\thetap = \mathrm{vec}(W)$ and $f_\thetap(x_{h,w}) = W^\top x_{h,w}$.
\end{assumption}

\begin{remark}
Assumption \ref{ass:linear} captures the canonical case used widely in
spectroscopy pipelines, including the FTIR and QCL tissue
classification work we draw motivation from. It is also the cleanest
case in which the Hessian of $\loss$ admits closed-form block analysis.
We discuss extensions to shallow nonlinear $f_\thetap$ in
Section~\ref{sec:discussion} and verify empirically in
Section~\ref{sec:experiments} that the theorem's scaling persists when
$f_\thetap$ is replaced by a small MLP.
\end{remark}

\begin{assumption}[Spatial model regularity]
\label{ass:reg}
$g_\phip$ is twice continuously differentiable in $\phip$ with
$L$-Lipschitz gradients.
\end{assumption}

\begin{assumption}[Sub-Gaussian inputs]
\label{ass:subg}
Per-pixel inputs $\{x_{h,w}\}$ are sub-Gaussian with second moment
$\Sigma = \E[x_{h,w} x_{h,w}^\top]$ of bounded condition number.
\end{assumption}

\begin{assumption}[Capacity gap]
\label{ass:gap}
$\Cg / \Cf \to \infty$.
\end{assumption}

Assumption~\ref{ass:gap} is the structural condition under which our
theorems become informative. Empirically it is satisfied with several
orders of magnitude to spare in modern spectral-spatial architectures
(in our motivating example, $\Cg/\Cf \approx 325$).

\begin{assumption}[Bounded input Jacobian of the classifier]
\label{ass:inputlip}
There exists a constant $\widetilde{L} > 0$ depending on the depth of
$g_\phip$, the bottleneck dimension $K$, the number of classes
$N_{\mathrm{cls}}$, and the data distribution --- but \emph{not} on
the characteristic width $M$ or the total spatial parameter count
$\Cg$ --- such that
\[
    \sup_{\phip \in \Phi_{\mathrm{reg}},\ Z \in \mathcal{Z}_{\mathrm{reg}}}
    \left\| \frac{\partial g_\phip(Z)}{\partial Z} \right\|_{\mathrm{op}}
    \;\leq\; \widetilde{L},
\]
uniformly over the training-time parameter regime
$\Phi_{\mathrm{reg}}$ and the induced feature domain
$\mathcal{Z}_{\mathrm{reg}}$.
\end{assumption}

\begin{remark}[Why Assumption~\ref{ass:inputlip} is separate from Assumption~\ref{ass:reg}]
Assumption~\ref{ass:reg} bounds the second derivatives of $g_\phip$
with respect to \emph{parameters}~$\phip$ (the standard
$L$-smoothness condition of the optimization literature).
Assumption~\ref{ass:inputlip} bounds the first derivative with respect
to \emph{input features} $Z$---a different object with a different
constant. Parameter-side smoothness does not imply input-side
Lipschitzness; we require both.

At random Kaiming/Xavier initialization, the operator norm of a
variance-scaled random weight matrix is $\Theta(1)$ in width $M$ by
the Bai--Yin law~\cite{bai1993limit} (see also
Vershynin~\cite{vershynin2018high}, Thm.~4.4.5;
Marchenko--Pastur~\cite{marchenko1967distribution} describes the
corresponding bulk spectrum), giving
$\|\partial g_\phip/\partial Z\|_{\mathrm{op}} = \Theta(1)$ in width
$M$ for standard CNN and MLP blocks. Under the maximal-update
parameterization ($\mu$P) of Yang and
Hu~\cite{yang2021tensor}, the bound persists throughout training as
$M \to \infty$. Under standard-parameterization training with weight
decay $\lambda > 0$ and bounded-norm iterates,
Assumption~\ref{ass:inputlip} follows from the boundedness of
$\|\phip\|$ combined with the per-layer product-of-operator-norms
bound; the constant $\widetilde{L}$ depends multiplicatively on the
depth $L$ but not on the width. Softmax self-attention is Lipschitz
only on bounded inputs~\cite{kim2021lipschitz}, so
$\mathcal{Z}_{\mathrm{reg}}$ should be taken compact whenever
$g_\phip$ uses attention---this is guaranteed by LayerNorm before
attention, which is standard practice.
\end{remark}

\subsection{Loss and gradients}

We train with per-pixel cross-entropy:
\begin{equation}
\label{eq:loss}
    \loss(\thetap, \phip) = \frac{1}{HW} \sum_{h=1}^{H} \sum_{w=1}^{W}
    \mathrm{CE}\!\left( F_{\thetap,\phip}(X)_{:,h,w},\ y_{h,w} \right),
\end{equation}
where $y_{h,w} \in \{1, \dots, N_{\mathrm{cls}}\}$ is the ground-truth
class. Cross-entropy plays an essential role in
Section~\ref{sec:thm2}: its gradient with respect to the logits
$\hat{y}$ is the prediction residual
$r = \mathrm{softmax}(\hat{y}) - y_{\mathrm{onehot}}$, which vanishes
exponentially as predictions saturate. We discuss other loss functions
sharing this vanishing-residual property in
Section~\ref{sec:discussion}.

The gradient of $\loss$ with respect to the spatial parameters is
computed by standard backpropagation. The gradient with respect to the
spectral parameters $\thetap$ admits the chain-rule decomposition
\begin{equation}
\label{eq:chain}
    \nabla_\thetap \loss
    = \underbrace{\frac{\partial \loss}{\partial \hat{y}}}_{\text{residual}}
    \cdot \underbrace{\frac{\partial \hat{y}}{\partial Z}}_{\text{spatial Jacobian}}
    \cdot \underbrace{\frac{\partial Z}{\partial \thetap}}_{\text{input}}.
\end{equation}
The three factors play distinct roles in our analysis. The residual is
small when predictions are confident. The spatial Jacobian depends on
the current state of $\phip$ and thus changes throughout training. The
final factor is the input itself and is stationary. The product
structure implies that if the residual decays to zero, the gradient to
$\thetap$ vanishes regardless of the other two factors. This
observation underlies Theorem~\ref{thm:twoscale}.

\subsection{The block Hessian}

The joint Hessian of $\loss$ decomposes into blocks aligned with the
two parameter groups:
\begin{equation}
\label{eq:block_hessian}
    H = \nabla^2 \loss(\thetap, \phip) =
    \begin{pmatrix}
        H_{\thetap\thetap} & H_{\thetap\phip} \\
        H_{\phip\thetap} & H_{\phip\phip}
    \end{pmatrix},
\end{equation}
with $H_{\thetap\thetap} \in \R^{\Cf \times \Cf}$,
$H_{\phip\phip} \in \R^{\Cg \times \Cg}$, and cross blocks of
appropriate sizes. Theorem~\ref{thm:hessian} bounds the condition
number of $H$ in terms of $\Cg / \Cf$ by separately bounding the
spectrum of each block.

For optimization-rate analysis it is often more tractable to work with
the Gauss-Newton approximation $G = J^\top J$, where $J = [J_\thetap,
J_\phip]$ is the Jacobian of the per-pixel residuals. The two matrices
coincide near interpolation, and Theorem~\ref{thm:hessian} is stated
for $G$.

\subsection{The Effective Gradient Ratio}

We introduce a real-time diagnostic that tracks the asymmetry in
gradient flow between the two modules.

\begin{definition}[Effective Gradient Ratio]
\label{def:egr}
At training step $t$ with parameters $(\thetap_t, \phip_t)$, the
Effective Gradient Ratio is
\[
    \EGR(t) \;:=\; \frac{\norm{\nabla_\thetap \loss(\thetap_t, \phip_t)}}
                        {\norm{\nabla_\phip \loss(\thetap_t, \phip_t)}}.
\]
\end{definition}

\begin{remark}
EGR is a unitless quantity directly computable inside any training
loop with negligible overhead. Section~\ref{sec:egr} establishes its
monotonicity along the slow-fast manifold (a consequence of
Theorem~\ref{thm:twoscale}) and proposes a threshold criterion for
detecting gradient starvation onset.
\end{remark}

\subsection{Notational conventions}

Throughout, $\norm{\cdot}$ denotes the Euclidean ($\ell_2$) norm for
vectors and the spectral norm for matrices; $\lambda_{\max}(\cdot)$ and
$\lambda_{\min}(\cdot)$ denote the largest and smallest eigenvalues of
a symmetric matrix; $\kappa(A) = \lambda_{\max}(A)/\lambda_{\min}(A)$
its condition number. We write $A \succeq B$ to mean $A - B$ is
positive semidefinite. The symbol $\Omega(\cdot)$ is used in the
standard asymptotic sense.

\begin{definition}[Module curvature disparity]
\label{def:dcurv}
The \emph{module curvature disparity} of the joint Gauss-Newton
matrix is
\[
    \mathcal{D}_{\mathrm{curv}}(G)
    \;:=\; \frac{\lambda_{\max}(G_{\phip\phip})}
                {\lambda_{\max}(G_{\thetap\thetap})}
    \;=\; \frac{\lambda_{\max}(J_{\phip}^\top J_{\phip})}
                {\lambda_{\max}(J_{\thetap}^\top J_{\thetap})},
\]
the ratio of the spatial Gauss-Newton block's top eigenvalue to the
spectral block's top eigenvalue.
\end{definition}

\begin{remark}[Why disparity, not condition number]
\label{rem:disparity_vs_kappa}
Theorem~\ref{thm:hessian} is a purely geometric statement about how
curvature concentrates between the two modules as the spatial network
widens. It is not a joint-Hessian condition-number bound: the joint
Gauss-Newton $G$ is rank-deficient (parameter directions orthogonal
to the classifier's range have exactly zero curvature), so
$\kappa(G) = \lambda_{\max}(G)/\lambda_{\min}(G)$ is trivially
infinite. The disparity $\mathcal{D}_{\mathrm{curv}}$ is
well-defined regardless of rank structure, uses only well-behaved
operator-norm quantities, and is exactly the ratio measured in
Experiment~1.1. Its growth with spatial capacity is the geometric
mechanism that motivates the two-timescale dynamical analysis of
Theorem~\ref{thm:twoscale}, but the dynamical claim requires its own
hypotheses and does not follow from $\mathcal{D}_{\mathrm{curv}}$
alone.
\end{remark}
